\scanpage{14}\mspage{13}
on trouve comme possibilités \unsure{optima|optimes} de definir la structure d'atelier sur $\unsure{\mathcal{F}|\mathfrak{F}}$ :
\[
\left\{\begin{aligned}
 &\text{Soit par}\quad\boxed{\leq,\quad\ll}\\
 &\text{soit par}\quad\leq,\quad\preccurlyeq\\
 &\text{soit par}\quad\ll,\quad\mathrel{\tikz[baseline=0pt,x=1em,y=1em]{\draw (-.38,.62)--(.38,.25)--(-.43,-.10); \draw (.34,.62)--(-.32,.25)--(.38,-.10);}}\ .
\end{aligned}\right.
\tag*{(1.28)}
\]
\margin{Non, \unsure{sa|ça} ne suffira pas, \unsure{d'apres|d'après} nouveau point de vue \ldots}

Dans la première mouture, j'ai choisi $\leq$, $\preccurlyeq$, en définissant en termes de ceux-ci
\[
\left\{\begin{aligned}
 &\ll\ \text{est la relation de préordre engendrée par}\\
 &\qquad\leq,\preccurlyeq,\quad \text{ou encore}\\
 &G\ll F\quad\Longleftrightarrow\quad\exists F'\ \text{avec}\quad G\leq F'\preccurlyeq F
\end{aligned}\right.
\tag*{(1.29)}
\]
% slanted
\margin{Je changerai le pt de vue page 20, et ne \unsure{supposerai} plus que $\ll$ soit déterminé ainsi en termes de $\leq,\preccurlyeq$ mais plutôt que $\preccurlyeq$ est déterminé par $\ll,\leq$}

Cela correspond au centre de gravité \deleted{conceptuel} \inserted{\unsure{géometrique|géométrique}} et intuitif, qui est bien sur la notion de subdivision, plutôt que sur celle de raffinement. Mais techniquement c'était un peu laborieux, de déduire les propriétés de transitivité et \inserted{surtout} de réflexivité pour $\ll$, cela amenait à mettre des axiomes en surnombre. Toujours techniquement, il semblerait que ce soit de partir de $\leq,\ll$ qui soit le plus simple. C'est l'approche que je vais suivre tantôt. Quant à la possibilité de passer par $\ll$, $\mathrel{\tikz[baseline=0pt,x=1em,y=1em]{\draw (-.38,.62)--(.38,.25)--(-.43,-.10); \draw (.34,.62)--(-.32,.25)--(.38,-.10);}}$, c'est plutôt une curiosité qu'autre chose,
